\section{Introduction}

At the core of many problems in optimisation, statistics, and machine learning lies the task of identifying a point that satisfies multiple constraints simultaneously~\cite[Ch. 8]{BOYDCONVEX}. When these constraints correspond to closed convex sets in a Hilbert space, this task takes the form of the \textit{best approximation problem} (BAP). Formally, we denote as $\mathcal{X}$ a real Hilbert space equipped with a norm $\| \cdot \|$. We consider a finite family of closed convex subsets $\{\mathcal{H}_i\}_{i=0}^{n-1}\subseteq\mathcal{X}$ with a non-empty intersection $\mathcal{H} := \bigcap_{i=0}^{n-1} \mathcal{H}_i$. Given an arbitrary point $x^\circ \in \mathcal{X}$, the solution to the BAP is the unique point $x^\star \in \mathcal{H}$ closest to $x^\circ$ with respect to the norm~\cite{BAUSCHKEBOOK}. The BAP is then formulated as the convex optimisation problem $x^\star = \Argmin_{x\in \mathcal{H}} \frac{1}{2} \|x - x^\circ\|^2$. This formulation is widely adopted in applications such as restricted least squares and isotonic regression in statistics, as well as in signal and image processing, where convex sets encode prior knowledge, like non-negativity or sparsity, in control engineering, where system constraints are typically modeled using convex sets, and in machine learning, where solutions to learning problems must often adhere to structural constraints~\cite{BAUSCHKESWISS}.

In this paper, we assume that $\mathcal{X}$ is a Euclidean space, $X \subseteq \Rset^p$, equipped with the $\ell_2$ norm $\twonorm{x}\eqdef (x_1^2+\dots+x_p^2)^{1/2}$, so that the BAP takes the form of a \textit{constrained quadratic program} (CQP):
\begin{align}\label{eq:projection}
x^\star := \mathcal{P}_{\mathcal{H}}(x^\circ) := \argmin_{x\in \bigcap_{i=0}^{n-1}\mathcal{H}_i}\frac{1}{2}\twonorm{x-x^\circ}^2.
\end{align}
Here, we further assume that $\mathcal{H}$ is a polyhedral set, $\mathcal{H}\eqdef\set{x\in\Rset^p}{A x \leq c}$, and that each $\mathcal{H}_i$ is a half-space,
\begin{align}\label{eq:sets_i}
\mathcal{H}_i\coloneqq\set{x\in\Rset^p}{\trans{f_i}x \leq c_i}, \quad i=1,\dots,n,
\end{align}
where the normal vector $f_i$ of each plane $H_i\coloneqq\set{x\in\Rset^p}{\trans{f_i}x = c_i}$ satisfies $\twonorm{f_i}=1$. While in some cases the projection onto $\mathcal{H}$ is easy to compute, for arbitrary polyhedral sets the projection operator $\mathcal{P}_{\mathcal{H}}$ onto the intersection set is intractable and no closed-form solution to~\eqref{eq:projection} is known. In these cases, the solution to~\eqref{eq:projection} can be obtained using standard optimisation software packages. To solve~\eqref{eq:projection}, most solvers require introducing an additional variable $z\coloneqq Ax$ (e.g., \cite{osqp} or~\cite{ocpb:16}) that is projected onto the set $\set{z\in\Rset^n}{z \leq c}$, for which an explicit solution exists~\cite{BOYDCONVEX}. However, for large $n$, this variable augmentation can reduce the computational efficiency, particularly in high-dimensional or time-sensitive applications when the projection is part of an overarching algorithm~\cite{OPTIMNESTEROV,KEMPF20206542}.

% I would entirely remove this. The problem is that we're not talking about this in the remainder of the paper. If you'd re-run the fast-gradient using our modified Dyktra and show that it's faster you would need this paragraph.
% As a recent example, Diamond Light Source (DLS), the UK's national synchrotron, is expected to see an upgrade to its electron beam stabilisation control system. Existing control methods (modal decomposition~\cite{HEATH}) are no longer sufficient: this is due to the increased number of sensors (from 172 to 252) and actuators (from 173 to 396), as well as the introduction of two different types of corrector magnets. A devised model predictive control (MPC) implementation~\cite{MPCDLSII}, while promising, faces high computational demands for high-frequency (100 kHz) control. To accommodate input rate constraints, the Fast Gradient method was proposed in~\cite{KEMPF20206542}, where a CQP problem needs to be solved at each timestep as part of the main loop, which is to be run less than \SI{10}{\micro\second} for feasibility. Evaluating the solution to the CQP using standard solvers would be too expensive for this high-latency requirement. This difficulty necessitates iterative methods that construct the solution to the CQP in linear time using only the individual projection operators $\mathcal{P}_{\mathcal{H}_i}$. Amongst these, Dykstra's method offers a promising solution due to its simplicity and convergence guarantees.

As an alternative that can be faster in practice~\cite{KEMPF20206542}, \textit{Dykstra's projection algorithm}~\cite{DYKSTRA} is an iterative method designed to solve~\eqref{eq:projection}. It is an extension of the \textit{method of alternating projections} (MAP), which, given an initial point $x^\circ$, finds a point in the intersection of $\mathcal{H}_i$. Dykstra's algorithm modifies MAP by introducing auxiliary variables to achieve strong convergence to the specific projection of $x^\circ$ onto the intersection of $\mathcal{H}_i$ (or general convex sets). Qualitatively, Dykstra's auxiliary variables remember the direction vector of projection; these are added to the primary variables before subsequent projections and, as a result, the primary iterates of Dykstra's method asymptotically converge to~\eqref{eq:projection}. Although, like standard optimisation packages, Dykstra's method also introduces additional variables to solve~\eqref{eq:projection}, these appear only in vector additions rather than matrix-vector multiplications. This reduces the computational load and could make Dykstra's method a highly valuable algorithm to deploy in real-time, high-throughput applications for finding fast (approximate) solutions to CQPs. However, despite its strong theoretical guarantees, Dykstra's algorithm is known to suffer from a \emph{stalling} problem for certain classes of initial conditions~\cite{DYKSTRASTALLING}. In such cases, the iterates of Dykstra's method remain unchanged for a number of iterations. The duration of the stalling period is not known a priori, and, given the choice of starting point, can be made arbitrarily long. In general, the algorithm cannot be run for a fixed number of iterations with a guarantee that the output will be closer to the projection than the initial point, which can hinder deployment of Dykstra's method in practical applications.

In this paper, we address the stalling problem of Dykstra's method for polyhedral sets. Exploiting simplifications that arise in the polyhedral case, we derive the exact duration of the stalling period once stalling is detected (Theorem~\ref{thm:nstall}). This duration enables a modified version of Dykstra's algorithm (Algorithm~\ref{alg:fastforward}) that fast-forwards through the stalling period, resulting in improved convergence properties. Crucially, this modification preserves all convergence guarantees of the original algorithm (Corollary~\ref{cor:convergence}). This development represents a step towards practical deployment of Dykstra's algorithm in real-time settings, where its low-overhead projection steps could offer significant performance advantages.

The structure of the paper is as follows. Section~\ref{sec:background} reviews MAP, Dykstra's algorithm, and narrows the analysis to polyhedral sets, introducing the stalling phenomenon. Section~\ref{sec:main_result} proposes the fast-forward modification, and Section~\ref{sec:example} presents supporting numerical experiments, which demonstrate the effectiveness of the solution. Finally, Section~\ref{sec:conclusion} concludes with a discussion of potential accelerations, extensions, and applications. 

We use standard notation throughout: $\trans{A}$ denotes the transpose of vector or matrix $A$, and $\twonorm{A}$ its $\ell_{2}$-norm. The interior of a set $\mathcal{X}$ is denoted as $\text{int}\mathcal{X}$. Function composition is written as $f \circ g$, and the modulo and ceiling operators are denoted by $[\cdot]$ and $\lceil \cdot \rceil$, respectively.